\section{常见的几何模型（一）}
\subsection{焦点三角形}
\clearpage
\subsection{离心率}
\clearpage
\subsection{定比点差法}
\clearpage
\subsection{Archimedes三角形}
\clearpage
\subsection{圆锥曲线的光学性质}
\begin{theorem}[椭圆的光学性质]
    
\end{theorem}
\begin{example}[2025八省联考18]
    已知椭圆 $C$ 的方程为\(\dfrac{x^2}{4}+\dfrac{y^2}{3}=1\)，左、右焦点分别为 $F_1(-1,0)$，$F_2(1,0)$。
  设 $M$ 是坐标平面上的动点，且线段 $F_1M$ 的垂直平分线与 $C$ 恰有一个公共点，证明 $M$ 的轨迹为圆，并求该圆的方程。
\end{example}
\begin{jie}
 设线段 $F_1M$ 的垂直平分线为 $l$，$l$ 与椭圆 $C$ 有唯一的公共点，即 $l$ 是椭圆 $C$ 的切线，记切点为 $Q$。由椭圆光学性质，$F_1$ 发出的光线经椭圆镜面反射至 $F_2$，即光线 $QF_2$ 是光线 $F_1Q$ 经镜面 $l$ 反射后的光线：
$M, F_1$ 关于直线 $l$ 对称，只需要说明$M, Q, F_2$ 三点共线，进而
\[
|MF_2| = |QM| + |QF_2| = |QF_1| + |QF_2| = 4,
\]
即 $M$ 的轨迹是以 $F_2$ 为半经，以$ 4$ 为半径的圆，其方程为 $(x-1)^2 + y^2 = 16$。

最后解决\(M,Q,F_2\)共线的问题：假设\(M,Q,F_2\)不共线，此时有$|MF_2| < |QM| + |QF_2| = |QF_1| + |QF_2| = 4$。
直接连接\(M, F_2\)交\(l\)于另一点\(Q'\)，由切线的性质，此\(Q'\)一定在椭圆之外，此时\(|MF_2|=|Q'F_2|+|Q'M|=|Q'F_1| + |Q'F_2|>4\)，产生矛盾！
故\(M,Q,F_2\)一定共线。
\end{jie}
由此可归纳出此一般性的命题：
\begin{proposition}
    
\end{proposition}
\begin{theorem}[双曲线的光学性质]
    
\end{theorem}
\begin{theorem}[抛物线的光学性质]
    
\end{theorem}
\begin{example}
 设抛物线 $C: x^{2}=2 y$ 的焦点为$ F$， 过抛物线$  C$ 上的点 $P$ 作 $ C $ 的切线交 
  $x$ 轴于点  $M $， 若$  \angle F P M=30^{\circ}$ , 求 
线段$F P$的长。
\end{example}

\clearpage

